\documentclass[12pt]{article}
\usepackage[utf8]{inputenc}
\usepackage[T1]{fontenc}
\usepackage{hyperref}
\title{Sungkai Language Resources for K-12 Curriculum Development}
\author{Compiled by Claude Sonnet 5.5 (claude-sonnet-5-5) for Andr\'as Kornai}
\date{}
\begin{document}
\maketitle

\section{Introduction}
The code suu (Glottocode sung1271) is a retired ISO 639-3 code. According to Glottolog it was merged into Lampung Api [ljp] with effect from 2008-01-14. Sungkai in this sense is a dialect community within the Lampung Api group in Sumatra, Indonesia. Joshua Project, citing Ethnologue, says the Sungkai use a dialect of Lampung Api. Do not confuse it with the unrelated Aslian language Sungkai (Glottolog sung1267), which is Austroasiatic. I found no resources specific to Sungkai (suu) as a separate language.

\section{Primary Educational Resources}
No Sungkai-specific primers, textbooks or school programmes were found in the search results.

\section{Linguistic and Reference Materials}
\begin{itemize}
\item Glottolog bookkeeping entry for Sungkai (Indonesia), retired and merged into Lampung Api: \url{https://glottolog.org/resource/languoid/id/sung1271}
\item Austronesian Basic Vocabulary Database, Lampung Api (Sungkai) wordlist: \url{https://abvd.eva.mpg.de/austronesian/language.php?id=608}
\end{itemize}

\section{Religious and Community Resources}
\begin{itemize}
\item Joshua Project profile of the Lampung Sungkai people group: \url{https://joshuaproject.net/people_groups/15124}
\end{itemize}

\section{Digital and Audio Resources}
None found.

\section{Curriculum Development Context}
Because the code is retired, the practical route is through the broader Lampung Api language and its resources, which should be searched under that name (ljp). I did not do that search here. Sungkai-specific materials would need to be sought from the local community. The ABVD wordlist is a small comparative list and not a teaching resource.

\end{document}